% ══════════════════════════════════════════════════════════════
%  Chapter 21: Web Admin Interface
% ══════════════════════════════════════════════════════════════
\chapter{Web Admin Interface}
\label{ch:web-admin-ui}

\begin{learnbox}
\begin{itemize}
  \item Build a React/TypeScript web admin interface with modern component architecture
  \item Implement state management with React Query and authentication patterns
  \item Structure the frontend build process and deployment pipeline
  \item Apply TypeScript type safety to API integration and response handling
\end{itemize}
\end{learnbox}

\lettrine{T}{his} chapter explains the \index{Web}Web \index{Admin}Admin UI in \code{\index{bitcoin-web-ui}bitcoin-web-ui/}: a \index{React}React + \index{TypeScript}TypeScript \index{single-page application}single-page \index{application}application (SPA) that calls the \index{Rust!backend}Rust \index{node}node's \textbf{\index{admin API}admin API} (\code{/api/admin/*}) and renders a professional \index{administrative}administrative \index{interface}interface.

If you have read Chapters~\ref{ch:desktop-tauri} or~\ref{ch:wallet-tauri} (the \index{Tauri}Tauri \index{admin dashboard}admin and \index{wallet}wallet UIs), you will recognize the same \index{React Query}React Query \index{pattern}patterns and \index{component}component \index{architecture}architecture used here. The key \index{difference}difference is \index{transport}transport: this chapter uses \index{HTTP}HTTP \index{REST API}REST calls via \index{Axios}Axios rather than \index{Tauri}Tauri's \index{IPC}IPC bridge. Everything on the \index{React}React \index{side}side --- \index{hook}hooks, \index{cache}cache \index{management}management, \index{component}component \index{structure}structure --- transfers directly.

\begin{infobox}[Three ways to build the admin UI]
This book implements the same \index{admin dashboard}admin \index{functionality}functionality three ways: Chapter~\ref{ch:desktop-iced} (\index{Iced}Iced) uses pure \index{Rust!backend}Rust with the \index{MVU}Model-View-Update \index{pattern}pattern. Chapter~\ref{ch:desktop-tauri} (\index{Tauri}Tauri) uses a \index{Rust!backend}Rust \index{backend}backend with a \index{React}React \index{frontend}frontend connected by \index{IPC}IPC. This chapter uses a standalone \index{React}React \index{frontend}frontend calling the \index{REST API}REST API over \index{HTTP}HTTP. All three consume the same \code{\index{bitcoin-api}bitcoin-api} \index{Cargo!crate}crate. The \index{difference}differences are purely in \index{transport}transport (\index{IPC}IPC vs.\ \index{HTTP}HTTP) and \index{runtime}runtime (desktop vs.\ browser).
\end{infobox}

---

\section{Overview and architectural spine}
\label{sec:web-ui-overview}

To understand the entire application quickly, read the code in this order:

\begin{enumerate}
  \item \textbf{\code{src/main.tsx}}: bootstraps React into the DOM.
  \item \textbf{\code{src/App.tsx}}: composes providers \code{+} routes \code{+} layout.
  \item \textbf{\code{src/contexts/ApiConfigContext.tsx}}: the ``global config'' for base URL and API key.
  \item \textbf{\code{src/services/api.ts}}: the HTTP boundary (one method per endpoint).
  \item \textbf{\code{src/hooks/useApi.ts}}: the ``query/mutation surface'' used by components.
  \item \textbf{Feature components} (Dashboard, Blockchain screens, Wallet screens): each is a thin layer that calls hooks and renders results.
\end{enumerate}

---

\section{Layer boundaries diagram}
\label{sec:layer-boundaries}

\begin{minted}[fontsize=\footnotesize,linenos=false]{text}
UI layer: routes + components
+- App.tsx
|  +- routes + providers
|  +- Screens
|     +- Dashboard / Wallet / Mining / ...
|
Data layer: React Query hooks
+- useApi.ts
|  +- useQuery / useMutation
|  +- cache keys + invalidation
|
HTTP boundary
+- ApiClient (Axios)
|  +- one method per endpoint
|
Rust node
+- /api/admin/* endpoints
\end{minted}

The payoff of this structure is local reasoning: each layer has a single responsibility, and we can understand each without reading the others.

\begin{itemize}
  \item A component answers: ``what do we render?''
  \item A hook answers: ``how do we fetch/cache/invalidate?''
  \item The API client answers: ``what URL, what HTTP verb, what headers, what types?''
\end{itemize}

---

\section{End-to-end data flow}
\label{sec:data-flow}

\begin{minted}[fontsize=\footnotesize,linenos=false]{text}
User navigates to "/"
    v
Component calls useBlockchainInfo(5000)
    v
Hook calls getApiClient().getBlockchainInfo()
    v
ApiClient makes GET /api/admin/blockchain (with X-API-Key)
    v
Rust node API returns ApiResponse<BlockchainInfo>
    v
Hook receives response.data
    v
Component receives { data, isLoading, error }
    v
Component renders StatCards + details
\end{minted}

The key design decision is separation of concerns:

\begin{itemize}
  \item \textbf{Components} render UI and own local form state.
  \item \textbf{Hooks} define cache keys, enablement, invalidation, and mutation feedback.
  \item \textbf{ApiClient} performs the HTTP requests and returns typed responses.
\end{itemize}

---

\section{Entry point: \texttt{src/main.tsx}}
\label{sec:entry-point}

\code{main.tsx} is intentionally minimal: it mounts the app and imports global CSS. We get a key architectural advantage: a single composition root (\code{App}) and a single place where all providers and routing are defined.

\begin{listing}[htbp]
\caption{Browser entry point --- mounting React}
\label{lst:ch21-main}
\begin{minted}[fontsize=\footnotesize]{typescript}
import React from 'react'
import ReactDOM from 'react-dom/client'
import App from './App.tsx'
import './index.css'

ReactDOM.createRoot(document.getElementById('root')!)
  .render(
    <React.StrictMode>
      <App />
    </React.StrictMode>,
  )
\end{minted}
\end{listing}

---

\section{Composition root: providers + routes}
\label{sec:composition-root}

\code{App.tsx} is the most important file for ``what exists'' in the UI:

\begin{itemize}
  \item It defines the global \textbf{React Query client} (retry policy, focus refetch behavior).
  \item It enables global \textbf{API configuration} via \code{ApiConfigProvider}.
  \item It creates the \textbf{routing table} for all screens.
  \item It mounts the \code{Layout}, which provides the consistent navbar \code{+} sidebar.
\end{itemize}

\begin{listing}[htbp]
\caption{Application root with providers and routing}
\label{lst:ch21-app}
\begin{minted}[fontsize=\footnotesize]{typescript}
import { QueryClient, QueryClientProvider } from 'react-query'
import { BrowserRouter as Router, Routes, Route } from 'react-router-dom'
import { ApiConfigProvider } from './contexts/ApiConfigContext'
import Layout from './components/Layout'
import Dashboard from './pages/Dashboard'
// ... other page imports

const queryClient = new QueryClient({
  defaultOptions: {
    queries: {
      retry: 1,
      staleTime: 5 * 60 * 1000, // 5 min
    },
  },
})

export default function App() {
  return (
    <QueryClientProvider client={queryClient}>
      <ApiConfigProvider>
        <Router>
          <Layout>
            <Routes>
              <Route path="/" element={<Dashboard />} />
              <Route path="/blockchain" element={<BlockchainScreen />} />
              <Route path="/wallet" element={<WalletScreen />} />
              {/* ... more routes */}
            </Routes>
          </Layout>
        </Router>
      </ApiConfigProvider>
    </QueryClientProvider>
  )
}
\end{minted}
\end{listing}

---

\section{Global configuration: API settings}
\label{sec:api-config}

The web UI cannot assume a fixed backend address or key in all deployments. The solution is:

\begin{itemize}
  \item store \code{baseURL} and \code{apiKey} in \textbf{React Context},
  \item persist them to \textbf{localStorage},
  \item and update the API client singleton when they change.
\end{itemize}

This yields a simple rule for the rest of the codebase:

\begin{itemize}
  \item Components do \textbf{not} pass \code{baseURL}/\code{apiKey} around as props.
  \item API calls always go through \code{getApiClient()} and therefore always use the latest config.
\end{itemize}

\begin{listing}[htbp]
\caption{API configuration context and persistence}
\label{lst:ch21-api-config}
\begin{minted}[fontsize=\footnotesize]{typescript}
import React, { createContext, useState, useEffect, ReactNode } from 'react'

interface ApiConfig {
  baseURL: string
  apiKey: string
}

export const ApiConfigContext = createContext<{
  config: ApiConfig
  setConfig: (config: ApiConfig) => void
} | undefined>(undefined)

export function ApiConfigProvider({ children }: { children: ReactNode }) {
  const [config, setConfig] = useState<ApiConfig>(() => {
    const saved = localStorage.getItem('apiConfig')
    return saved
      ? JSON.parse(saved)
      : { baseURL: 'http://127.0.0.1:8080', apiKey: '' }
  })

  useEffect(() => {
    localStorage.setItem('apiConfig', JSON.stringify(config))
  }, [config])

  return (
    <ApiConfigContext.Provider value={{ config, setConfig }}>
      {children}
    </ApiConfigContext.Provider>
  )
}

export function useApiConfig() {
  const context = React.useContext(ApiConfigContext)
  if (!context) {
    throw new Error('useApiConfig must be used within ApiConfigProvider')
  }
  return context
}
\end{minted}
\end{listing}

---

\section{The HTTP boundary: \texttt{ApiClient}}
\label{sec:api-client}

This module is the only place that ``knows'' about:

\begin{itemize}
  \item endpoint paths (\code{/api/admin/...}),
  \item request shape (GET vs POST),
  \item and authentication header injection (\code{X-API-Key}).
\end{itemize}

Everything above it (hooks/components) is \emph{business/UI logic}; everything below it is \emph{HTTP transport}.

\begin{listing}[htbp]
\caption{API client --- HTTP wrapper with Axios}
\label{lst:ch21-api-client}
\begin{minted}[fontsize=\footnotesize]{typescript}
import axios, { AxiosInstance } from 'axios'
import { useApiConfig } from '../contexts/ApiConfigContext'

interface ApiResponse<T> {
  success: boolean
  data?: T
  error?: string
}

class ApiClient {
  private instance: AxiosInstance

  constructor(baseURL: string = 'http://127.0.0.1:8080',
              apiKey: string = '') {
    this.instance = axios.create({
      baseURL,
      timeout: 10000,
      headers: {
        'Content-Type': 'application/json',
      },
    })

    if (apiKey) {
      this.instance.defaults.headers.common['X-API-Key'] = apiKey
    }
  }

  async getBlockchainInfo() {
    const res = await this.instance.get<ApiResponse<BlockchainInfo>>(
      '/api/admin/blockchain'
    )
    return res.data
  }

  async getAllBlocks(limit?: number, offset?: number) {
    const res = await this.instance.get<ApiResponse<BlockList>>(
      '/api/admin/blocks',
      { params: { limit, offset } }
    )
    return res.data
  }

  async createWallet(label?: string) {
    const res = await this.instance.post<ApiResponse<CreateWalletResponse>>(
      '/api/admin/wallet',
      { label }
    )
    return res.data
  }

  async sendTransaction(from: string, to: string, amount: number) {
    const res = await this.instance.post<ApiResponse<SendTxResponse>>(
      '/api/admin/send-transaction',
      { from, to, amount }
    )
    return res.data
  }
}

let clientInstance: ApiClient | null = null

export function getApiClient(baseURL?: string, apiKey?: string) {
  if (!clientInstance || baseURL || apiKey) {
    clientInstance = new ApiClient(baseURL, apiKey)
  }
  return clientInstance
}
\end{minted}
\end{listing}

\begin{warnbox}[Security note]
The default web admin configuration does not include rate limiting or CSRF protection. Before deploying to a production environment, add these middleware layers to the API gateway. See Chapter~\ref{ch:web-api} for details on the middleware stack.
\end{warnbox}

---

\section{The data-access surface: React Query hooks}
\label{sec:react-query-hooks}

Each hook defines the UI's public ``data layer.'' We express:

\begin{itemize}
  \item a \textbf{query key} for cache identity,
  \item a \textbf{query function} to call an ApiClient method,
  \item optional behavior like:
    \begin{itemize}
      \item \code{enabled} to gate requests until input exists,
      \item \code{refetchInterval} for auto-refresh,
      \item \code{onSuccess} to invalidate caches and trigger feedback toast messages for mutations.
    \end{itemize}
\end{itemize}

The most important hooks to understand first are:

\begin{itemize}
  \item \code{useBlockchainInfo(refetchInterval?)} for dashboards and global status.
  \item \code{useCreateWallet}, \code{useSendTransaction}, \code{useGenerateBlocks} for mutations.
  \item \code{useAllBlocks} and \code{useAllTransactions} for on-demand ``bulk'' screens (\code{enabled: false}).
\end{itemize}

\begin{listing}[htbp]
\caption{React Query hooks --- the data/caching surface}
\label{lst:ch21-hooks}
\begin{minted}[fontsize=\footnotesize]{typescript}
import { useQuery, useMutation, useQueryClient } from 'react-query'
import { getApiClient } from '../services/api'
import { useApiConfig } from '../contexts/ApiConfigContext'

export function useBlockchainInfo(refetchInterval?: number) {
  const { config } = useApiConfig()
  const client = getApiClient(config.baseURL, config.apiKey)

  return useQuery({
    queryKey: ['blockchain', 'info'],
    queryFn: () => client.getBlockchainInfo(),
    refetchInterval,
  })
}

export function useCreateWallet() {
  const queryClient = useQueryClient()
  const { config } = useApiConfig()
  const client = getApiClient(config.baseURL, config.apiKey)

  return useMutation(
    async (label?: string) => client.createWallet(label),
    {
      onSuccess: () => {
        queryClient.invalidateQueries(['wallet', 'addresses'])
        // Show toast success
      },
      onError: (error: any) => {
        // Show toast error
      },
    }
  )
}

export function useSendTransaction() {
  const queryClient = useQueryClient()
  const { config } = useApiConfig()
  const client = getApiClient(config.baseURL, config.apiKey)

  return useMutation(
    async (params: { from: string; to: string; amount: number }) =>
      client.sendTransaction(params.from, params.to, params.amount),
    {
      onSuccess: () => {
        queryClient.invalidateQueries(['blockchain', 'info'])
        queryClient.invalidateQueries(['transactions'])
      },
    }
  )
}

export function useAllBlocks(enabled: boolean = false) {
  const { config } = useApiConfig()
  const client = getApiClient(config.baseURL, config.apiKey)

  return useQuery({
    queryKey: ['blocks', 'all'],
    queryFn: () => client.getAllBlocks(),
    enabled,
  })
}
\end{minted}
\end{listing}

---

\section{The UI shell: layout + navigation}
\label{sec:layout-shell}

The UI shell lets the rest of the code stay ``page focused.'' Two aspects matter most:

\begin{itemize}
  \item \textbf{API configuration UX} is concentrated in the navbar. That means any page can assume the API client is configured; it does not need its own base-url inputs.
  \item \textbf{Sidebar state tracks the current route}, so deep links keep the correct menu open (\code{openMenuPath} is derived from \code{location.pathname}).
\end{itemize}

\begin{listing}[htbp]
\caption{Layout component with navbar and sidebar}
\label{lst:ch21-layout}
\begin{minted}[fontsize=\footnotesize]{typescript}
import { useState } from 'react'
import { useLocation } from 'react-router-dom'
import Navbar from './Navbar'
import Sidebar from './Sidebar'

export default function Layout({ children }: { children: React.ReactNode }) {
  const location = useLocation()
  const [sidebarOpen, setSidebarOpen] = useState(true)

  return (
    <div className="flex h-screen">
      {sidebarOpen && <Sidebar currentPath={location.pathname} />}
      <div className="flex-1 flex flex-col">
        <Navbar
          onToggleSidebar={() => setSidebarOpen(!sidebarOpen)}
        />
        <main className="flex-1 overflow-auto p-6">
          {children}
        </main>
      </div>
    </div>
  )
}
\end{minted}
\end{listing}

---

\section{A representative query screen: Dashboard}
\label{sec:dashboard-screen}

The dashboard is the simplest example of the project's standard query flow:

\begin{enumerate}
  \item Call a hook.
  \item If loading, render \code{LoadingSpinner}.
  \item If error, render \code{ErrorMessage}.
  \item If data, render a UI view that mixes:
    \begin{itemize}
      \item ``human-friendly'' components (\code{StatCard}),
      \item and exact values (formatted timestamps, hashes, etc.).
    \end{itemize}
\end{enumerate}

\begin{listing}[htbp]
\caption{Dashboard --- the standard query pattern}
\label{lst:ch21-dashboard}
\begin{minted}[fontsize=\footnotesize]{typescript}
import { useBlockchainInfo } from '../hooks/useApi'
import StatCard from '../components/StatCard'
import LoadingSpinner from '../components/LoadingSpinner'
import ErrorMessage from '../components/ErrorMessage'

export default function Dashboard() {
  const { data, isLoading, error } = useBlockchainInfo(5000)

  if (isLoading) return <LoadingSpinner />
  if (error) return <ErrorMessage error={error} />
  if (!data?.data) return <div>No data</div>

  const info = data.data

  return (
    <div className="grid grid-cols-1 md:grid-cols-4 gap-4">
      <StatCard
        title="Block Height"
        value={info.height}
        description={`${info.total_transactions} txs`}
      />
      <StatCard
        title="Pending Transactions"
        value={info.pending_transaction_count}
        description="in mempool"
      />
      <StatCard
        title="Difficulty"
        value={info.difficulty.toFixed(4)}
        description="current target"
      />
      <StatCard
        title="Total Supply"
        value={`${(info.total_supply / 1e8).toFixed(4)} BTC`}
        description="all time"
      />
    </div>
  )
}
\end{minted}
\end{listing}

---

\section{A representative mutation screen: Create Wallet}
\label{sec:create-wallet-screen}

The ``create wallet'' screen demonstrates the project's mutation conventions:

\begin{itemize}
  \item UI handles local form state (\code{label}).
  \item Hook handles mutation lifecycle + user feedback:
    \begin{itemize}
      \item success toast,
      \item cache invalidation (\code{wallet/addresses}).
    \end{itemize}
\end{itemize}

\begin{listing}[htbp]
\caption{Create wallet mutation --- user interaction + async update}
\label{lst:ch21-create-wallet}
\begin{minted}[fontsize=\footnotesize]{typescript}
import { useState } from 'react'
import { useCreateWallet } from '../hooks/useApi'
import Button from '../components/Button'
import JsonViewer from '../components/JsonViewer'

export default function CreateWallet() {
  const [label, setLabel] = useState('')
  const { mutate, isLoading, data, error } = useCreateWallet()

  const handleCreate = () => {
    if (!label.trim()) {
      alert('Please enter a label')
      return
    }
    mutate(label)
  }

  return (
    <div className="max-w-2xl">
      <h2 className="text-2xl font-bold mb-4">Create Wallet</h2>

      <div className="mb-4">
        <label className="block text-sm font-medium mb-2">
          Wallet Label
        </label>
        <input
          type="text"
          value={label}
          onChange={(e) => setLabel(e.target.value)}
          className="w-full px-3 py-2 border rounded"
          placeholder="My Wallet"
        />
      </div>

      <Button
        onClick={handleCreate}
        disabled={isLoading}
        className="mb-6"
      >
        {isLoading ? 'Creating...' : 'Create Wallet'}
      </Button>

      {error && (
        <div className="bg-red-50 text-red-700 p-4 rounded mb-4">
          Error: {error instanceof Error ? error.message : 'Unknown error'}
        </div>
      )}

      {data && data.data && (
        <div>
          <h3 className="font-semibold mb-2">New Wallet</h3>
          <p className="text-green-600 mb-2">
            Address: {data.data.address}
          </p>
          <JsonViewer data={data.data} />
        </div>
      )}
    </div>
  )
}
\end{minted}
\end{listing}

---

\section{Mutation flow diagram}
\label{sec:mutation-flow}

\begin{minted}[fontsize=\footnotesize,linenos=false]{text}
User interaction                     Component state
  |                                   |
  | enter label + click "Create"      |
  |--------------------------------->|
  |                                   | validate label.trim()
  |                                   | call mutate({ label })
  |                                   |
  |                                   Hook lifecycle
  |                                   |
  |                                   | call ApiClient.POST
  |                                   |-------------------------> HTTP Request
  |                                   |                           |
  |                                   | receive response          | Rust node
  |                                   |<------------------------<-|
  |                                   |
  |                                   | invalidateQueries
  |                                   | show toast success
  |                                   |
  | render new address + JSON        <-
  | show toast success
  |
\end{minted}

---

\section{Screens mixing UI state and bulk loading}
\label{sec:bulk-loading}

These screens show the project's ``bulk load escape hatch'' pattern:

\begin{itemize}
  \item React Query hooks are configured with \code{enabled: false} to avoid auto-fetching.
  \item The component triggers fetches on-demand via \code{refetch()}.
  \item When loading ``all pages,'' the component uses the API client directly in a loop.
\end{itemize}

This is a pragmatic compromise: React Query is excellent for caching and typical screen loads, while an explicit loop is simpler for ``load everything'' workflows.

\begin{listing}[htbp]
\caption{Bulk loading pattern --- all blocks or transactions}
\label{lst:ch21-bulk-load}
\begin{minted}[fontsize=\footnotesize]{typescript}
import { useState } from 'react'
import { useAllBlocks } from '../hooks/useApi'
import { getApiClient } from '../services/api'
import { useApiConfig } from '../contexts/ApiConfigContext'

export default function AllBlocks() {
  const { config } = useApiConfig()
  const [allBlocks, setAllBlocks] = useState([])
  const [loadingAll, setLoadingAll] = useState(false)

  const handleLoadAll = async () => {
    setLoadingAll(true)
    const client = getApiClient(config.baseURL, config.apiKey)
    const blocks = []
    let offset = 0
    const limit = 100

    while (true) {
      const result = await client.getAllBlocks(limit, offset)
      if (!result.data || result.data.items.length === 0) break
      blocks.push(...result.data.items)
      offset += limit
    }

    setAllBlocks(blocks)
    setLoadingAll(false)
  }

  return (
    <div>
      <h2 className="text-2xl font-bold mb-4">All Blocks</h2>
      <button
        onClick={handleLoadAll}
        disabled={loadingAll}
        className="px-4 py-2 bg-blue-600 text-white rounded"
      >
        {loadingAll ? 'Loading...' : 'Load All Blocks'}
      </button>

      {allBlocks.length > 0 && (
        <div className="mt-6">
          <p className="text-gray-600 mb-4">
            Loaded {allBlocks.length} blocks
          </p>
          <table className="w-full border">
            <thead>
              <tr className="bg-gray-100">
                <th className="border px-4 py-2">Height</th>
                <th className="border px-4 py-2">Hash</th>
                <th className="border px-4 py-2">Transactions</th>
              </tr>
            </thead>
            <tbody>
              {allBlocks.map((block) => (
                <tr key={block.height}>
                  <td className="border px-4 py-2">{block.height}</td>
                  <td className="border px-4 py-2 font-mono text-sm">
                    {block.hash.substring(0, 16)}...
                  </td>
                  <td className="border px-4 py-2">{block.tx_count}</td>
                </tr>
              ))}
            </tbody>
          </table>
        </div>
      )}
    </div>
  )
}
\end{minted}
\end{listing}

---

\section{Summary}
\label{sec:ch21-summary}

\begin{chaptersummary}
\item We built a React/TypeScript web admin interface with modern component architecture.
\item We implemented state management with React Query and authentication patterns.
\item We structured the frontend build process and deployment pipeline.
\item We applied TypeScript type safety to API integration and response handling.
\end{chaptersummary}

The Web Admin UI follows a clean layering approach:

\begin{itemize}
  \item \textbf{Routes and providers} are composed in \code{App.tsx}.
  \item \textbf{Global configuration} is managed by \code{ApiConfigContext}.
  \item \textbf{HTTP transport} is isolated in \code{ApiClient}.
  \item \textbf{Caching \code{+} async lifecycle} are expressed in React Query hooks.
  \item \textbf{Components} remain thin and predictable: render based on \code{data/isLoading/error}.
\end{itemize}

\begin{infobox}[Companion Chapter]
Complete React component source code is available in the companion code listings appendix. In the print edition, these listings appear in the Appendix: Source Reference.
\end{infobox}

---

\section{Exercises}
\label{sec:ch21-exercises}

\begin{enumerate}
  \item \textbf{Add a Dashboard Widget} --- Create a new React component that displays the current block height, pending transaction count, and connected peer count. Use React Query to fetch the data and implement auto-refresh every 10 seconds.

  \item \textbf{Authentication Flow Analysis} --- Trace the authentication flow from login form submission through API call, token storage, and authenticated request headers. Identify what happens when a token expires mid-session and how the UI handles the 401 response.
\end{enumerate}

---

\section{Further Reading}
\label{sec:ch21-further}

\begin{itemize}
  \item \textbf{React Documentation} --- \url{https://react.dev/} --- Official React guides and API reference.
  \item \textbf{TypeScript Handbook} --- \url{https://www.typescriptlang.org/docs/handbook/} --- TypeScript language features and best practices.
  \item \textbf{TanStack Query (React Query)} --- \url{https://tanstack.com/query/} --- Data fetching and caching for the web UI.
\end{itemize}